\section{2d/3d 对偶：CS/WZW 与 AdS$_3$/CFT$_2$ 是两码事}
\label{sec:duality}

用户原话：“这里的对偶跟 AdS$_3$/CFT$_2$ 有什么区别？CS/WZW 对偶和全息对偶是两码事吗，这个得搞清楚。”

答案是：\textbf{是两码事}，而且不是同一种“对偶”。

\subsection{两种关系的类型不同}

\begin{center}
\small
\begin{tabular}{@{}p{1.9cm}p{5.0cm}p{5.0cm}@{}}
\toprule
维度 & CS/WZW（本文，1988） & AdS$_3$/CFT$_2$（1997 后） \\
\midrule
体理论 & 三维 Chern--Simons：无度规（作用量 (1.3) 不含 $g$）\src{354}，态空间 $\cH_\Sigma$ 有限维 \src{366--367} & 三维量子引力（含度规，渐近 AdS$_3$ 边界条件）\\
边界理论 & 二维 WZW／手征代数；其共形块空间（原文 p.370 以 “conformal blocks” 名之）\src{370} & 二维 CFT（一般大 $c$）\\
关系类型 & 量子化态空间 = 共形块空间：\textbf{态空间等同} \src{368--370}；三维侧只在 $X\times S^1$ 以 $Z=\dim\cH_X$ 接触二维数据 \src{387--388} & 体配分函数 = 边界配分函数：\textbf{动力学等价} \\
体自由度 & 无度规自由度；$\cH_\Sigma$ 有限维 \src{366--367} & 有：引力自由度、BTZ 黑洞、黑洞熵 \\
引力子 & 无（理论里没有度规） & 有 \\
中心荷 & $c=k\dim G/(k+h^\vee)$ \src{390 (4.41)}；$k\to\infty$ 时 $c\to\dim G$ \src{361}；$k\in\mathbb Z$ \src{354} & Virasoro 渐近对称，$c=3l/2G$（Brown--Henneaux），连续\\
可证性 & 可逐条检验（辛商、Narasimhan--Seshadri、BWB） & 无定理；只有弱耦合/对称性/可积性检验 \\
全息字典 & 不存在“体场↔边界算子”字典 & 存在（GKPW 型）\\
\bottomrule
\end{tabular}
\end{center}

\subsection{唯一的接触点，以及为什么它不是等同}

还要说清一件事：本文的“三维—二维接触”只有一条通道——在 $X\times S^1$ 型流形上，三维配分函数等于 $X$ 上 Hilbert 空间的维数，即 $Z(X\times S^1)=\dim\cH_X$ \src{387--388}。\textbf{没有“体配分函数 = 边界配分函数”这样的等式}；这正是它不能叫全息对偶的技术原因。

确有一个接触点：\textbf{三维引力可以写成 Chern--Simons 理论}（规范群取 $ISO(2,1)$ 或其覆盖群）。Witten 在本文只一句带过，细节在姊妹篇 \src{397}（文献 [38]，预印本 HEP-88/32）。

但这个改写\textbf{不是}等同：它要求\textbf{非紧}规范群，要求特定边界条件；本文的 CS 用\textbf{紧}群、没有边界引力自由度。改写保留的是经典运动方程，并没有把三维引力变成“有限维态空间”的拓扑场论。

\begin{keypoint}
\textbf{原文从未出现 AdS、从未出现“全息”这个词}（1988 年还没有这个说法）。本文 §5 结尾只说“正是这类考虑促成了对 2+1 维引力的研究” \src{397}——这是动机链，不是等价性陈述。
\end{keypoint}

\subsection{为什么容易混}

两个原因。第一，两者都叫“3d/2d 对应”，字面上都是“三维理论对应二维理论”。第二，CS/WZW 的边界理论确实是二维 CFT，而 AdS$_3$/CFT$_2$ 的边界也是二维 CFT——但前者的“二维理论”是\textbf{态空间的组织方式}（共形块），后者的“二维理论”是\textbf{另一个动力学理论}。

一句话区分：\textbf{CS/WZW 是同一套量子化数据的两种写法；AdS$_3$/CFT$_2$ 是两个理论之间的动力学等价。}

\begin{boundary}
\begin{itemize}[leftmargin=1.4em]
\item 已确证：本文的对应是态空间级 \src{368--369,394--397}；原文无 AdS/全息字样；姊妹篇 [38] 是 3d 引力=CS 的出处。
\item 我们补的：对照表的六个维度（体自由度、引力子、中心荷来源、可证性、字典）。
\item 未核（后世文献卷页）：Brown--Henneaux 的 $c=3l/2G$、Maldacena 1997、GKPW 字典；Achúcarro--Townsend 1986。
\end{itemize}
\end{boundary}
